\documentclass[a4paper,11pt]{article}
\usepackage[utf8]{inputenc}
\usepackage{geometry}
\geometry{margin=2.5cm}

\usepackage{hyperref}
\hypersetup{colorlinks=true,linkcolor=blue,urlcolor=blue}

\usepackage{setspace}
\onehalfspacing

\begin{document}

\begin{flushright}
Attila Német\\
29 rue Jean-Pierre Sauvage\\
2514 Luxembourg\\
+352-661-118922\\
\href{mailto:attila.nemet@gmail.com}{attila.nemet@gmail.com}
\end{flushright}

\vspace{1cm}

\begin{flushleft}
Digital Learning Hub\\
Service de la Formation des Adultes\\
Ministère de l'Éducation nationale,\\
de l'Enfance et de la Jeunesse\\
Esch-sur-Alzette
\end{flushleft}

\vspace{0.5cm}

\begin{flushright}
Luxembourg, le 19 mars 2026
\end{flushright}

\vspace{0.5cm}

\noindent\textbf{Objet : Candidature au poste d'Accompagnateur socio-éducatif 42 Luxembourg (réf. E00041327) (réf. F00041329) (réf. M00041332)}

\vspace{0.5cm}

\noindent Madame, Monsieur,

\vspace{0.3cm}

\noindent C'est avec un grand enthousiasme que je vous adresse ma candidature pour le poste d'accompagnateur socio-éducatif au sein de 42 Luxembourg. En tant qu'apprenant actif à 42 Luxembourg, je vis au quotidien la réalité du campus, la pédagogie par les pairs et les défis auxquels font face les apprenants. Cette perspective unique, combinée à plus de vingt ans d'expérience professionnelle internationale, me permet de proposer un accompagnement à la fois crédible, empathique et ancré dans la réalité de l'école.

\vspace{0.3cm}

\noindent \textbf{Une connaissance intime de 42 Luxembourg.} Je connais personnellement les Piscines, le stress qu'elles génèrent, les moments de doute et de découragement, mais aussi la satisfaction immense de franchir chaque étape. J'ai vécu l'apprentissage par les pairs, les évaluations croisées et le travail collaboratif qui font l'identité de 42. Cette expérience me permet de comprendre en profondeur ce que traversent les apprenants et de leur offrir un espace d'écoute authentique et bienveillant.

\vspace{0.3cm}

\noindent \textbf{Un parcours multiculturel au service de l'inclusion.} Au cours de ma carrière chez Ericsson, j'ai travaillé dans plus de 20 pays -- du Japon à l'Irlande, de l'Égypte à la Suède -- au sein d'équipes extrêmement diversifiées. Cette expérience m'a appris l'écoute active, la diplomatie, la médiation interculturelle et le respect des différences. J'ai régulièrement accompagné et soutenu des collègues juniors dans des environnements à forte pression, en développant une approche bienveillante et orientée vers la solution.

\vspace{0.3cm}

\noindent \textbf{Des compétences en communication et en gestion de situations complexes.} En tant que Technical Lead à Dublin et Business Analyst à Budapest, j'ai coordonné des équipes pluridisciplinaires, géré des situations de stress élevé lors de déploiements critiques en réseau opérationnel, et assuré la liaison entre les équipes techniques et managériales. Ces expériences m'ont doté d'une capacité éprouvée à communiquer avec clarté, à désamorcer les tensions et à maintenir un climat de travail constructif.

\vspace{0.3cm}

\noindent \textbf{Une compréhension du marché IT et de l'insertion professionnelle.} Mon parcours de près de trente ans dans le secteur des technologies de l'information me donne une vision concrète du marché de l'emploi IT, des compétences recherchées et des réalités du monde professionnel. Je suis en mesure d'orienter les apprenants dans leur préparation à l'insertion professionnelle et de les mettre en lien avec les bons interlocuteurs.

\vspace{0.3cm}

\noindent \textbf{Un profil multilingue.} Je maîtrise l'anglais (C2) et le français (B2), et je parle également l'allemand (B1), le roumain (C2) et le hongrois (langue maternelle), ce qui me permet d'interagir avec un large éventail d'apprenants dans un environnement aussi multiculturel que celui du Luxembourg.

\vspace{0.3cm}

\noindent Titulaire d'un Master de l'Université du Luxembourg (niveau 7), je suis convaincu que mon profil atypique -- à la croisée de l'expérience technique, de l'engagement humain et de la connaissance directe de 42 -- constitue un atout pour accompagner les apprenants dans leur parcours de formation et les aider à s'épanouir tant sur le plan personnel que professionnel.

\vspace{0.3cm}

\noindent Je me tiens à votre entière disposition pour un entretien au cours duquel je pourrai vous présenter plus en détail ma motivation et mon projet pour ce poste.

\vspace{0.5cm}

\noindent Je vous prie d'agréer, Madame, Monsieur, l'expression de mes salutations distinguées.

\vspace{1.5cm}

\noindent Attila Német

\end{document}
